\section{The occurrence-indexed ledger}
\label{sec:ledger}

Fix a universe of smooth projective varieties over a characteristic-zero
field and choose a skeleton.  This harmless choice only avoids set-theoretic
duplication; isomorphisms remain among the generating correspondences.

\begin{definition}[Occurrence provider]
\label{def:provider}
An occurrence provider assigns to every smooth projective $X$ a finite set
$A_X$ and a function $m_X:A_X\to\N_{>0}$.  Its expanded occurrence set is
\[
 \Occ(X):=\{(a,j):a\in A_X,\ 1\le j\le m_X(a)\}.
 \tag{2.1}
\]
The provider includes bijections, compatible with isomorphisms, of the
following forms:
\begin{align}
 \Occ(X\sqcup Y)&\simeq \Occ(X)\sqcup\Occ(Y),
                                                        \tag{2.2}\label{eq:disjoint}\\
 \Occ(\Bl_ZX)&\simeq
   \Occ(X)\sqcup\bigl(\Occ(Z)\times\{1,\ldots,c-1\}\bigr),
                                                        \tag{2.3}\label{eq:blowup}\\
 \Occ(\PP_X(E))&\simeq
   \Occ(X)\times\{0,\ldots,r-1\},
                                                        \tag{2.4}\label{eq:projective}
\end{align}
where $Z\subset X$ is smooth of codimension $c\ge2$ and $E$ has rank
$r\ge2$.
The labels $j$ have no structure: a provider bijection is used only up to
permuting the occurrences above a fixed element of $A_X$.
\end{definition}

Let \(\mathscr O\) be the disjoint union of the sets \(\Occ(X)\).  First
identify $(a,j)$ with $(a,j')$ for all labels over the same local block.
Then use isomorphisms and each provider bijection, in either direction, to
generate an equivalence relation on \(\mathscr O\).  Let
\(\Pi\) be the associated thin groupoid: its objects are \(\mathscr O\), and
there is a unique arrow $x\to y$ exactly when the two occurrences are
equivalent.  Put
\[
 H:=\pi_0(\Pi),
 \qquad
 \mathcal L:=\Sym^{\sqcup}(H)=\N^{(H)}.
 \tag{2.5}\label{eq:ledger}
\]
Thus \(\mathcal L\) is the free commutative monoid on the occurrence classes.
Write $q:\mathscr O\to H$ for the quotient map and define
\[
 \mathcal B(X):=\sum_{x\in\Occ(X)}[q(x)]\in\mathcal L.
 \tag{2.6}\label{eq:block-ledger}
\]

The blowup correspondence gives the commutative square
\begin{equation}
\begin{CD}
 \Occ(\Bl_ZX) @>{\sim}>>
 \Occ(X)\sqcup\bigl(\Occ(Z)\times\{1,\ldots,c-1\}\bigr)\\
 @V{q}VV @VV{q^{\sqcup}}V\\
 H @= H .
\end{CD}
\tag{2.7}\label{diag:occurrence-descent}
\end{equation}
It records descent on objects.  It does not assert an isomorphism between
unrelated representatives of a common class.

\begin{theorem}[Universal marker ledger]
\label{thm:ledger}
The assignment \(\mathcal B\) satisfies
\begin{align}
 \mathcal B(X\sqcup Y)&=\mathcal B(X)+\mathcal B(Y),
                                                        \tag{2.8}\label{eq:ledger-disjoint}\\
 \mathcal B(\Bl_ZX)&=\mathcal B(X)+(c-1)\mathcal B(Z),
                                                        \tag{2.9}\label{eq:ledger-blowup}\\
 \mathcal B(\PP_X(E))&=r\mathcal B(X).                 \tag{2.10}\label{eq:ledger-projective}
\end{align}
For every commutative monoid $M$, a function $w_0:\mathscr O\to M$ that
is constant on the generating correspondences factors uniquely as
\[
 \mathscr O\xrightarrow{q}H\xrightarrow{w}M.
\]
The function $w$ in turn extends uniquely to a monoid homomorphism
\(\overline w:\mathcal L\to M\), and
\(I_w(X):=\overline w(\mathcal B(X))\) satisfies the folded versions of
\eqref{eq:ledger-disjoint}--\eqref{eq:ledger-projective}.
\end{theorem}

\begin{proof}
The provider bijections match every occurrence, including multiplicity.
Changing a copy index is one of the generating arrows, so the construction is
independent of the temporary labelling.  After applying $q$, summing the two sides of
\eqref{eq:disjoint}--\eqref{eq:projective} gives
\eqref{eq:ledger-disjoint}--\eqref{eq:ledger-projective}.  A function on
\(\mathscr O\) is constant on the generated equivalence relation precisely
when it factors uniquely through $H$.  The universal property of the free
commutative monoid then gives the unique extension \(\overline w\).  Applying
that homomorphism to the three ledger identities proves the last assertion.
\end{proof}

This formulation makes multiplicity bookkeeping literal.  If a component has
degree $m$, it contributes $m$ objects to \(\Occ(X)\) and hence $m$
copies of one generator to \(\mathcal B(X)\).  No symmetric-algebra convention
or passage to a Grothendieck group is needed.

\subsection{The effective dimension quotient}

For $e\ge0$, let $H_{\le e}\subset H$ be the classes represented by an
occurrence on some smooth projective variety of dimension at most $e$.  The
projection
\[
 \rho_e:\mathcal L=\N^{(H)}\longrightarrow
 \mathcal L_{>e}:=\N^{(H\setminus H_{\le e})}
 \tag{2.11}\label{eq:dimension-quotient}
\]
deletes the low-dimensional coordinates.  Because both monoids are free and
effective, \(\rho_e(\ell)\ne0\) exactly when \(\ell\) has a positive
coefficient on a class outside $H_{\le e}$.

\begin{theorem}[Weak-factorization quotient]
\label{thm:birational-quotient}
Let $d\ge2$.  The assignment
\[
 X\longmapsto \rho_{d-2}(\mathcal B(X))
 \tag{2.12}\label{eq:birational-ledger}
\]
is a birational invariant of smooth projective $d$-folds.  In particular,
if \(\rho_{d-2}(\mathcal B(X))\ne0\), then $X$ is irrational.
\end{theorem}

\begin{proof}
Projective weak factorization expresses a birational map between smooth
projective $d$-folds as blowups and blowdowns along smooth centers of
codimension at least two \cite[Theorem~0.1.1]{AKMW}.  Every center has
dimension at most $d-2$, so \(\rho_{d-2}(\mathcal B(Z))=0\).
Equation~\eqref{eq:ledger-blowup} therefore shows that
\eqref{eq:birational-ledger} is unchanged at each step.

For projective space, \eqref{eq:ledger-projective} applied over a point gives
\(\mathcal B(\PP^d)=(d+1)\mathcal B(\operatorname{Spec}k)\).  The point
classes lie in $H_{\le d-2}$, since $d\ge2$, so the quotient ledger of
\(\PP^d\) is zero.  A rational $X$ has the same quotient ledger, proving the
contrapositive.
\end{proof}
